\documentclass[11pt]{article}
\usepackage[margin=1in]{geometry}
\usepackage{amsmath,amssymb,amsthm}
\usepackage{array}
\usepackage[round]{natbib}
\usepackage{hyperref}
\usepackage{microtype}

\newtheorem{theorem}{Theorem}
\newtheorem{proposition}{Proposition}
\newtheorem{lemma}{Lemma}
\newtheorem{corollary}{Corollary}
\newtheorem{definition}{Definition}
\theoremstyle{remark}
\newtheorem{remark}{Remark}

\newcommand{\R}{\mathbb{R}}
\newcommand{\E}{\mathbb{E}}
\newcommand{\Prob}{\mathbb{P}}
\newcommand{\ind}[1]{\mathbf 1\{#1\}}
\newcommand{\tv}[1]{\lVert #1\rVert_{\mathrm{TV}}}

\title{Herding as the Contragredient of Conformal Prediction}
\author{Peter Cotton\\\texttt{peter.cotton@microprediction.com}}
\date{July 2026}

\begin{document}
\maketitle

\begin{abstract}
The contragredient of a group action is the same action carried by the dual space: transform the
functionals instead of the points. We show that this is the exact relation between conformal
calibration, whose validity proof transports an evaluation functional over fixed scores, and
herding, which transports points against fixed functionals; both control one argument of the
bilinear pairing $\langle h,\nu\rangle=\int h\,d\nu$ while a cumulative discrepancy stays bounded
by a constant, which is where both $O(1/n)$ rates come from. The relation is useful exactly
insofar as proofs cross it, and we cross it in both directions through a stability lemma for the
pairing: the coverage floor of \citet{barberpananjady2026} for split conformal on $\beta$-mixing
series factors into the lemma's two Lipschitz legs --- buffer deletion as a measure perturbation,
the switch coefficient as a distributional total variation --- which is why their rate is linear
in $\tau/n$ where blocking arguments give $\sqrt{\tau/n}$; and, in the reverse direction, the
herding argument transfers to the calibration set, so that a low-discrepancy design makes
realized coverage deterministic to $O(1/n)$. The conformal mechanism (late 1990s) predates
herding (2009); the contragredient is herding.
\end{abstract}

\section{Introduction}

Two finite-sample constructions, developed independently, obtain $O(1/n)$ rates where sampling
arguments give $O(n^{-1/2})$: the calibration step of conformal prediction, whose mechanism dates
to the confidence machines of the late 1990s (see \citealp{vovk2005}, for the canonical
treatment, and \citealp{lei2018}, for the split form used here), and herding
\citep{welling2009}.

The relation between them is named in the title, so we define the name at once. Given a group
acting on a vector space, the \emph{contragredient} is the induced action on the dual space: if
points transform by $\rho(\pi)$, functionals transform by $\rho^{*}(\pi)=\rho(\pi^{-1})^{\top}$,
and the pairing of a functional with a point is preserved. Herding acts on points: it selects
$x_1,x_2,\dots$ so that fixed test functionals average correctly. Conformal calibration acts on
functionals: the scores are never moved, the acceptance threshold is steered, and the validity
proof transports the functional ``evaluate the test coordinate'' across the sample. One mechanism
is the other with the action moved to the dual side.

A relation of this kind earns its keep only if arguments cross it. Sections~\ref{sec:conformal}
and~\ref{sec:herding} give the two mechanisms parallel formal statements;
Section~\ref{sec:duality} states the duality on the pairing
$\langle h,\nu\rangle=\int h\,d\nu$; Section~\ref{sec:transfer} crosses it in both directions,
factoring the dependent-data coverage floor of \citet{barberpananjady2026} through a stability
lemma for the pairing and transferring the herding argument to the design of calibration sets.
Remarks, including the relation to the Bayesian-quadrature reading of \citet{snellgriffiths2025},
close in Section~\ref{sec:remarks}.

\section{Conformal calibration}\label{sec:conformal}

Fix a miscoverage level $\alpha\in(0,1)$ and write $p=1-\alpha$. Let $S_1,\dots,S_n$ be
calibration scores and $S_{n+1}$ a test score, assumed almost surely distinct, and let
$k=\lceil p\,(n+1)\rceil$, assumed $\le n$. Split conformal prediction declares coverage when
$S_{n+1}\le q_n$, where $q_n=S_{(k)}$ is the $k$-th smallest calibration score.

The choice of $q_n$ is a feedback rule on the acceptance test $h_q(s)=\ind{s\le q}$.

\begin{lemma}[Steering identity]\label{lem:steer}
With $q_n=S_{(k)}$ and the scores distinct,
\[
\sum_{i=1}^{n}\bigl(h_{q_n}(S_i)-p\bigr)\;=\;k-pn\;\in\;[\,p,\;1+p\,).
\]
In particular the cumulative acceptance-rate discrepancy is bounded by an absolute constant
independent of $n$, and the empirical acceptance rate differs from $p$ by at most $(1+p)/n$.
\end{lemma}

\begin{proof}
Exactly $k$ of the $n$ calibration scores satisfy $S_i\le S_{(k)}$, so the sum equals $k-pn$.
Since $p(n+1)\le k<p(n+1)+1$, we have $p\le k-pn<1+p$.
\end{proof}

The control-theoretic reading of conformal prediction is established in the online setting: the
level update of adaptive conformal inference \citep{gibbscandes2021} is integral feedback on
realized miscoverage, and \citet{angelopoulos2023pid} make the identification explicit as
proportional--integral--derivative control. There the controlled quantity is the time-average
miscoverage along a sequence; Lemma~\ref{lem:steer} is the batch counterpart, a single dead-beat
correction applied to the calibration set before any test data arrive.

Validity rests on a deterministic counting identity concerning the \emph{placement} of the test
index.

\begin{proposition}[Counting identity]\label{prop:count}
Let $s_1,\dots,s_{n+1}$ be distinct reals and, for $i\in\{1,\dots,n+1\}$, call placement $i$
\emph{covered} if $s_i$ is at most the $k$-th smallest of $\{s_j\}_{j\ne i}$. Then placement $i$
is covered if and only if the rank of $s_i$ among all $n+1$ values is at most $k$, and
consequently exactly $k$ of the $n+1$ placements are covered.
\end{proposition}

\begin{proof}
Placement $i$ is covered iff at most $k-1$ of the other $n$ values are smaller than $s_i$, iff the
rank of $s_i$ among all $n+1$ values is at most $k$. The ranks are a permutation of
$\{1,\dots,n+1\}$, so exactly $k$ indices qualify.
\end{proof}

\begin{corollary}[Marginal validity]\label{cor:valid}
If $(S_1,\dots,S_{n+1})$ is exchangeable with almost surely distinct components, then
$\Prob\{S_{n+1}\le q_n\}=k/(n+1)\ge 1-\alpha$.
\end{corollary}

\begin{proof}
The coverage event $\{S_{n+1}\le q_n\}$ is the event that placement $n+1$ is covered, since
$q_n$ is the $k$-th smallest of the values other than $S_{n+1}$. By exchangeability the indicator
of ``placement $i$ is covered'' has the same expectation for every $i$; average
Proposition~\ref{prop:count} over $i$.
\end{proof}

No probabilistic inequality enters before the final step, and the final step is an identity in
distribution rather than a concentration bound.

\section{Herding}\label{sec:herding}

Let $\phi:\mathcal X\to\mathcal H$ be a feature map into a Hilbert space and
$\mu=\E_{P}\,\phi(X)$ a target mean. Herding \citep{welling2009} generates points by the recursion
\begin{equation}\label{eq:herding}
x_{t}\;=\;\arg\max_{x\in\mathcal X}\,\langle w_{t-1},\,\phi(x)\rangle,
\qquad
w_{t}\;=\;w_{t-1}+\mu-\phi(x_{t}),
\end{equation}
with $w_0$ given.

\begin{proposition}[Herding discrepancy]\label{prop:herd}
For any sequence produced by \eqref{eq:herding},
\[
\Bigl\|\frac1n\sum_{t=1}^{n}\phi(x_t)-\mu\Bigr\|\;=\;\frac{\|w_0-w_n\|}{n}.
\]
In particular, if $\sup_t\|w_t\|\le B$ then the feature-mean discrepancy is at most $2B/n$.
\end{proposition}

\begin{proof}
Telescoping \eqref{eq:herding} gives $w_n=w_0+n\mu-\sum_{t\le n}\phi(x_t)$; rearrange and divide
by $n$.
\end{proof}

Conditions under which the weights remain bounded are given by \citet{welling2009} in finite
dimension and analysed in general via the identification of \eqref{eq:herding} with the
Frank--Wolfe algorithm on the marginal polytope \citep{bach2012}; \citet{chen2010} treat the
kernel case. Under such conditions herding attains $O(1/n)$ where i.i.d.\ sampling attains
$O(n^{-1/2})$.

\section{The duality}\label{sec:duality}

Sections~\ref{sec:conformal} and~\ref{sec:herding} exhibit the same ledger: a cumulative
discrepancy --- $k-pn$ in Lemma~\ref{lem:steer}, $w_0-w_n$ in Proposition~\ref{prop:herd} ---
bounded by a constant and divided by $n$. To compare the mechanisms, write both in terms of the
bilinear pairing between bounded test functions $h$ and finite measures $\nu$,
\[
\langle h,\nu\rangle\;=\;\int h\,d\nu .
\]
Each mechanism controls one argument of this pairing and holds the other fixed. Herding fixes the
family of tests (the coordinates of $\phi$, whose $\nu$-averages are to match $\mu$) and
constructs the measure, $\nu_n=\frac1n\sum_t\delta_{x_t}$. Conformal calibration fixes the
measure --- the empirical distribution of scores --- and constructs the test $h_{q_n}$; its
validity proof additionally moves the evaluation index, as follows.

Let $\rho$ denote the permutation representation of the symmetric group $\mathfrak S_{n+1}$ on
$\R^{n+1}$, $\rho(\pi)e_i=e_{\pi(i)}$, with contragredient
$\rho^{*}(\pi)=\rho(\pi^{-1})^{\top}$, so that $\rho^{*}(\pi)e^{*}_{j}=e^{*}_{\pi(j)}$ and
$\langle\rho^{*}(\pi)f,\rho(\pi)v\rangle=\langle f,v\rangle$.

\begin{proposition}[Two descriptions of the orbit average]\label{prop:dual}
For $s\in\R^{n+1}$ with distinct entries let
$\chi_i(s)=\ind{s_i\le Q_k(s_{-i})}$, where $Q_k(s_{-i})$ is the $k$-th smallest entry of $s$ with
coordinate $i$ deleted. Then for every $\pi\in\mathfrak S_{n+1}$,
\[
\chi_{n+1}\bigl(\rho(\pi)s\bigr)\;=\;\chi_{\pi^{-1}(n+1)}(s).
\]
Consequently the average of $\chi_{n+1}$ over rearrangements of the data equals the average of
$s\mapsto\chi_i(s)$ over placements $i$ of the evaluation index --- equivalently, over the orbit
$\{\rho^{*}(\pi)\,e^{*}_{n+1}\}$ of the evaluation functional at fixed data --- and
Corollary~\ref{cor:valid} may be proved on either side of the pairing.
\end{proposition}

\begin{proof}
$Q_k$ is invariant under permutations of its $n$ arguments, and coordinate $n+1$ of $\rho(\pi)s$
is $s_{\pi^{-1}(n+1)}$; substitute.
\end{proof}

\begin{remark}[Direction, and what the name does]\label{rmk:direction}
Permutation representations are orthogonal, so $\rho^{*}$ is equivalent to $\rho$; the
contragredient adds no new representation, and its content here is the location of the action ---
on the covector side of the pairing for conformal, on the vector side for herding. Chronology
fixes the direction of the phrase: the conformal mechanism predates herding by roughly a decade,
so it is herding that realises the contragredient, in the primal space of generated points, with
active selection of $x_t$ replacing the transport of an evaluation index. The dependence
analysis of \citet{barberpananjady2026} is more recent than either; the duality claims kinship of
proofs, not priority of authors. Whether the relation is more than a name is settled by
Section~\ref{sec:transfer}, where arguments cross it in both directions.
\end{remark}

The correspondence is summarised below.

\begin{center}
\begin{tabular}{r@{\qquad}l@{\qquad}l}
 & \textbf{Conformal calibration} & \textbf{Herding} \\[2pt]
controlled side & test $h_{q_n}$ (and evaluation index) & measure $\nu_n$ \\
fixed side & empirical score measure & feature tests $\phi$ \\
ledger & $k-pn\in[p,1+p)$ & $w_0-w_n$, $\|w_t\|\le B$ \\
rate & $(1+p)/n$ & $2B/n$ \\
probabilistic input & exchangeability & none (deterministic target) \\
\end{tabular}
\end{center}

\begin{remark}[An image]
In a gambling idiom: herding moves the balls to balance a fixed wheel; conformal prediction
leaves the balls where they fell and relabels the slots under the pointer.
\end{remark}

\section{Moving arguments across the pairing}\label{sec:transfer}

The vehicle is the joint Lipschitz property of the pairing.

\begin{lemma}[Pairing stability]\label{lem:pair}
For tests $h,h'$ and finite measures $\nu,\nu'$,
\[
\bigl|\langle h,\nu\rangle-\langle h',\nu'\rangle\bigr|
\;\le\;\|h-h'\|_\infty\,\nu'(\mathbf 1)\;+\;\|h\|_\infty\,\tv{\nu-\nu'},
\]
where $\tv{\cdot}$ is the total-variation norm. Perturbations of either argument cost linearly,
and nothing concentrates.
\end{lemma}

\begin{proof}
Add and subtract $\langle h,\nu'\rangle$ and bound each difference by H\"older.
\end{proof}

\subsection{Toward conformal: the Barber--Pananjady floor}\label{sec:bp}

We first state the result exactly. For a score vector $S=(S_1,\dots,S_{n+1})$ and parameters
$k,\tau$, \citet{barberpananjady2026} define two subvectors of length $n+1-\tau$:
$\Delta^{0}_{k,\tau}(S)$ deletes a block of $\tau$ entries while keeping $S_{n+1}$ in the final
slot, and $\Delta^{1}_{k,\tau}(S)$ deletes the same number while rotating $S_k$ into the final
slot. Their \emph{switch coefficient} is the distributional distance
\[
\Psi_{k,\tau}(S)\;=\;d_{\mathrm{TV}}\bigl(\mathcal L(\Delta^{0}_{k,\tau}(S)),\,
\mathcal L(\Delta^{1}_{k,\tau}(S))\bigr),
\qquad
\bar\Psi_\tau\;=\;\frac1{n+1}\sum_{k=1}^{n+1}\Psi_{k,\tau}(S).
\]
For a pretrained, memoryless score their Theorem~5 gives the floor
\begin{equation}\label{eq:bp}
\Prob\{S_{n+1}\le q_n\}\;\ge\;1-\alpha-\min_{0\le\tau\le n}
\Bigl\{\tfrac{\tau}{n+1}+\bar\Psi_\tau\Bigr\},
\end{equation}
their Corollary~6 bounds $\bar\Psi_\tau\le 2\beta(\tau)$ for stationary $\beta$-mixing sequences,
their Theorem~8 supplies a matching ceiling, and their Theorem~7 shows \eqref{eq:bp} is tight
over this class up to the factor $(1-\alpha)/4$. Predictors with memory $L$ and scores trained on
an initial block are handled with the same structure at the cost of extra buffer terms.

The two steps of their argument instantiate the two legs of Lemma~\ref{lem:pair}. The deletion
step perturbs the measure argument: removing $\tau$ of $n+1$ points moves the empirical score
measure by at most $2\tau/(n+1)$ in total-variation norm, and the image of that perturbation
under the quantile functional is the level shift $\tau/(n+1)$ in \eqref{eq:bp} --- the
insertion--deletion stability of order statistics is the pairing's measure-side Lipschitz leg,
specialised to $h$ an indicator. The switch step perturbs on the test side in expectation: the
evaluation functional is transported from slot $n+1$ to slot $k$, and the price is, by the
definition of $\Psi_{k,\tau}$, a total variation between the laws of the two arrangements ---
the lemma applied to the distribution of the pairing rather than to a fixed realisation. Both
legs are linear, so the floor is linear in $\tau/n$; a blocking argument
\citep{oliveira2024} bounds the same quantity through an empirical-process maximal inequality
and pays the Monte Carlo rate $\sqrt{\tau/n}$. Linearity is not an accident of their proof; it is
the linearity of the pairing.

\subsection{Toward herding: designed calibration}\label{sec:design}

In the reverse direction the herding argument itself crosses the pairing: apply the bounded-%
discrepancy mechanism not to the generation of points for integration but to the selection of a
calibration set.

\begin{proposition}[Low-discrepancy calibration]\label{prop:design}
Let $\bar F$ be a continuous score distribution and suppose the calibration values
$S_1,\dots,S_n$ (distinct, obtained by any design --- herding, kernel thinning, a quantile grid)
satisfy $\sup_s|\hat F_n(s)-\bar F(s)|\le D$, where $\hat F_n$ is their empirical distribution
function. Then the realized coverage of the split threshold against $\bar F$ satisfies,
deterministically,
\[
\bigl|\bar F(S_{(k)})-\tfrac{k}{n+1}\bigr|\;\le\;D+\tfrac1n .
\]
\end{proposition}

\begin{proof}
$\hat F_n(S_{(k)})=k/n$, so $|\bar F(S_{(k)})-k/n|\le D$, and $k/n-k/(n+1)=k/(n(n+1))\le 1/n$.
\end{proof}

A quantile grid achieves $D\le 1/(2n)$; herding and kernel thinning achieve $D=O(1/n)$ up to
logarithms in more general settings \citep{chen2010,dwivedi2021}. Realized coverage is then
within $O(1/n)$ of nominal on \emph{every} draw: the dispersion of realized coverage collapses
from the sampling order $n^{-1/2}$ to the herding order $n^{-1}$. The price is stated by the
counting argument itself: selection by score value breaks the exchangeability of the calibration
set with the test point, so the symmetry proof of Corollary~\ref{cor:valid} no longer applies and
the guarantee is the deterministic one above, against the marginal law. The trade --- symmetry
for determinism --- is the design of a measure-side controller for an index-side guarantee.

\section{Remarks}\label{sec:remarks}

\begin{remark}[Quadrature]
\citet{snellgriffiths2025} recast the conformal computation as Bayesian quadrature over the
calibration scores. That is a measure-side reading of the same pairing: quadrature weights on
scores on one side, a transported evaluation functional on the other.
\end{remark}

\begin{remark}[Limits]
The duality concerns the finite-sample ledger only. The passage from
Proposition~\ref{prop:count} to Corollary~\ref{cor:valid} is exchangeability's, or, in
Section~\ref{sec:bp}, its mixing surrogate's; no bookkeeping supplies it. Nor does control of
thresholds and labels create information: it re-levels coverage while conditional sharpness
remains a property of the scores \citep{cotton_marginally}. Whether the steering can be aimed
conditionally --- selecting or reweighting calibration slots to balance the rank count given the
input, as herding aims its particles at a target --- is open; localized and weighted conformal
methods are the nearest existing steps, and the same information-theoretic ceiling applies to
them.
\end{remark}

\bibliographystyle{plainnat}
\bibliography{../notes}

\end{document}
